%%%%%%%%%%%%%%%%%%%%%%%%%%%%%%%%%%%%%%%%%%%%%%%%%%%%%%%%%%%%%%%%%%%%%%%%%%%%%%
%  Unit III, second half: E-mail Marketing (Chapters 4-7)
%%%%%%%%%%%%%%%%%%%%%%%%%%%%%%%%%%%%%%%%%%%%%%%%%%%%%%%%%%%%%%%%%%%%%%%%%%%%%%

\chapter{Effective E-mail Campaigns}

\syllabusstrip{E-mail Marketing: Effective E-mail Campaigns.}

\begin{learningoutcomes}
\begin{itemize}
\item What e-mail marketing is, and why it is an owned channel.
\item E-mail 1.0 against E-mail 2.0.
\item The advantages, and the liabilities that come with them.
\item Opt-in against spam, and the legal position.
\item The anatomy of an effective e-mail --- three rules and a checklist.
\item The content arsenal: nine e-mail types in three strategic groups.
\item Transactional against promotional e-mail.
\end{itemize}
\end{learningoutcomes}

\section{What E-mail Marketing Is}

\begin{definitionbox}[E-mail Marketing]
\keytermas{E-mail marketing}{e-mail marketing} is the use of e-mail to deliver
commercial and relationship-building messages to a permission-based list of
subscribers, in order to acquire, engage, convert and retain customers.
\end{definitionbox}

It is a vital tool used throughout the entire sales funnel\ix{sales funnel} --- engaging
prospects, maintaining interest and driving key revenue for online retailers at a
fraction of the cost of direct mail.

\begin{keybox}[Why E-mail Is Structurally Different]
Every other channel rents its audience. Social reach is granted by an algorithm
that can withdraw it; search ranking is granted by an algorithm that can change.
An e-mail list is an \textbf{owned asset}: no intermediary stands between the
message and the recipient, which is why e-mail has the highest conversion and
retention rate of any digital channel.
\end{keybox}

\section{E-mail 1.0 Against E-mail 2.0}

\begin{mmlongtable}{The generational shift in e-mail marketing}{tab:email12}%
{@{}P{0.44\textwidth}P{0.49\textwidth}@{}}
\rowcolor{ocre}\thd{E-mail 1.0} & \thd{E-mail 2.0}\\
\midrule
\endfirsthead
\rowcolor{ocre}\thd{E-mail 1.0} & \thd{E-mail 2.0}\\
\midrule
\endhead
Simple one-way messages & Behaviour-based triggers\\
\rowcolor{tablealt}
Static auto-responders & Segmented targeting based on purchase history and
interests\\
Mass ``batch and blast'' distribution & Dynamic lifecycle automation\\
\rowcolor{tablealt}
Untracked engagement & Deep tracking metrics\\
\end{mmlongtable}

\section{Advantages}

\begin{exambox}[Four Advantages Over Traditional Methods --- 2 Marks]
\pyq{CIE-II, 26 May 2026 --- 2 M}
Low cost; global reach; easy performance tracking; personalised communication.
\end{exambox}

\begin{itemize}
\item \sfb{Cost and speed} --- cheaper and faster than traditional direct mail.
\item \sfb{Tracking} --- real-time visibility into opens, clicks and conversions.
\item \sfb{Personalisation} --- dynamic tailoring based on collected data.
\item \sfb{Global reach} at no incremental cost per recipient.
\item \sfb{Ownership of the channel} --- unlike social media, the list is an
      owned asset that no algorithm can throttle.
\item \sfb{Highest conversion and retention rate} of any digital channel.
\item \sfb{Direct and permission-based} --- the recipient asked to hear from you.
\end{itemize}

\section{Disadvantages and Liabilities}

\begin{itemize}
\item \sfb{Spam and legal risk} --- filters block messages, and anti-spam
      legislation carries severe penalties.
\item \sfb{Engagement decay} --- constant tweaking is required to prevent list
      fatigue.
\item \sfb{Code rendering} --- appearance breaks across varying e-mail clients,
      servers and settings, for example when images are blocked by default.
\item \sfb{Hidden costs} --- scaling campaigns requires design, management and
      software investment, potentially upward of \$1,000 per month.
\item \sfb{Deliverability management} --- sender reputation, authentication (SPF,
      DKIM, DMARC) and list hygiene require ongoing attention.
\item \sfb{Inbox competition} --- the average inbox is crowded, so open rates are
      structurally limited.
\end{itemize}

\section{Opt-In Against Spam}

\begin{warnbox}[The Legal Point]
The CAN-SPAM Act levies fines of up to \textbf{\$16,000 per violation} for
non-compliance. Permission is not a courtesy; it is a legal requirement. Indian
and EU equivalents --- the IT Act rules and GDPR --- impose comparable
obligations.
\end{warnbox}

\begin{mmlongtable}{Spam against opt-in}{tab:optin}%
{@{}P{0.17\textwidth}P{0.38\textwidth}P{0.38\textwidth}@{}}
\rowcolor{ocre}\thd{} & \thd{Spam} & \thd{Opt-In}\\
\midrule
\endfirsthead
\rowcolor{ocre}\thd{} & \thd{Spam} & \thd{Opt-In}\\
\midrule
\endhead
\rlab{Acquisition method} & Unsolicited; default web checkboxes; purchased lists &
Willingly provided via lead forms in exchange for incentives\\
\rowcolor{tablealt}
\rlab{Expectation setting} & Blind blasting; unpredictable frequency & Clear
promises about e-mail type --- sales against tips --- and frequency\\
\rlab{Brand impact} & Degrades reputation; perceived as a nuisance & Builds
long-term trust; establishes authority; boosts upsells\\
\rowcolor{tablealt}
\rlab{List quality} & High volume, low intent & Micro-focused, highly targeted
engagement\\
\end{mmlongtable}

\begin{definitionbox}[Single and Double Opt-In]
\keytermas{Single opt-in}{single opt-in} adds the address immediately on
submission. \keytermas{Double opt-in}{double opt-in} requires the subscriber to
confirm through a verification e-mail --- producing a smaller but far cleaner and
more engaged list, improving deliverability and providing documentary proof of
consent.
\end{definitionbox}

\begin{exambox}[Why Segmentation Matters in E-mail --- 2 Marks]
\pyq{SEE, Jun/Jul 2025 --- 2 M}
Segmentation divides the list by demographics, purchase history, behaviour,
engagement level or lifecycle stage so that each group receives relevant content.
It raises open and click rates, reduces unsubscribes and spam complaints,
improves deliverability by maintaining engagement, and enables personalisation
--- which is the single largest driver of e-mail conversion.
\end{exambox}

\section{Anatomy of an Effective E-mail}

\subsection{The three rules}

\begin{frameworkbox}[The Three Rules]
\textbf{1. One topic.} Restrict the e-mail to a single marketing message.
Multiple messages create a ``busy'' user experience; focus drives higher revenue
per message.

\textbf{2. Attractive, simple design.} The design must complement the message,
not distract from it. Design on the assumption that mobile clients may disable
images.

\textbf{3. Clear call to action\ix{call to action}.} The reader must never wonder what to do next;
the path to purchase must be frictionless.
\end{frameworkbox}

\begin{keybox}[The Result of Experience]
Continuously A/B test designs and offers in order to learn customer behaviour. No
template is permanently correct.
\end{keybox}

\subsection{Checklist for an effective marketing e-mail}

\begin{itemize}
\item A \sfb{subject line} that is specific, benefit-led and under about 50
      characters --- this alone determines the open rate.
\item A \sfb{recognisable, consistent sender name} --- people open based on who
      sent it as much as on what it says.
\item \sfb{Preheader text} that extends the subject line rather than repeating
      it.
\item \sfb{Personalisation beyond the first name} --- reference actual behaviour
      or past purchases.
\item \sfb{Above-the-fold value} --- the offer must be visible without scrolling.
\item A \sfb{single, prominent, repeated CTA button} with active verb text.
\item \sfb{Responsive, mobile-first design} --- the majority of opens are on
      mobile.
\item \sfb{Alt text on every image}, so the message survives image blocking.
\item A \sfb{plain-text fallback} version.
\item A \sfb{visible unsubscribe link} and the sender's physical address --- both
      legally required.
\item \sfb{Tested rendering} across major clients before sending.
\end{itemize}

\section{The Content Arsenal --- Nine Types of Marketing E-mail}

\begin{keybox}[Name the Group, Not Just the Type]
The nine templates fall into three strategic groups. Naming the group and the
rules for each type is what separates a full-mark answer from a list.
\end{keybox}

\begin{figure}[htbp]
\centering
\begin{tikzpicture}[font=\sffamily\footnotesize,
    grp/.style={rounded corners=4pt,line width=1pt,align=center,
                text width=3.35cm,inner sep=7pt,font=\sffamily\scriptsize,
                text=slate}]
  \node[grp,draw=warnred,fill=warnpale] at (0,0)
    {\textbf{\textcolor{warnred}{GROUP 1}}\\ \textbf{DRIVE REVENUE}\\[5pt]
     Promotional\\ New Inventory\\ Reorder};
  \node[grp,draw=ocre,fill=ocrepale] at (4.05,0)
    {\textbf{\textcolor{ocredark}{GROUP 2}}\\ \textbf{BUILD RELATIONSHIPS}\\[5pt]
     Newsletter\\ Welcome\\ Educational\\ Product Advice};
  \node[grp,draw=deepblue,fill=deepbluepale] at (8.1,0)
    {\textbf{\textcolor{deepblue}{GROUP 3}}\\ \textbf{GATHER INTELLIGENCE}\\[5pt]
     Testimonials\\ Surveys};
\end{tikzpicture}
\caption[The nine e-mail types in three strategic groups]%
        {The content arsenal. The group determines the tone and the metric: Group~1
         is judged on revenue, Group~2 on engagement, Group~3 on response rate.}
\label{fig:ninetypes}
\end{figure}

\subsection{Group 1 --- drive revenue}

\begin{mmlongtable}{Group 1: revenue e-mails}{tab:group1}%
{@{}P{0.15\textwidth}P{0.28\textwidth}P{0.50\textwidth}@{}}
\rowcolor{ocre}\thd{Type} & \thd{Goal} & \thd{Rules}\\
\midrule
\endfirsthead
\rowcolor{ocre}\thd{Type} & \thd{Goal} & \thd{Rules}\\
\midrule
\endhead
\rlab{Promotional} & Entice immediate purchase & Keep it short; make the offer
front and centre; create a sense of urgency with active language\\
\rowcolor{tablealt}
\rlab{New Inventory} & Alert customers to new arrivals & Send immediately upon
arrival; prioritise a strong product photograph; minimise text; tease the reveal
in the subject line\\
\rlab{Reorder} & Replenish consumables & Make the buying process simple with a
``Reorder Now'' button; remind them of the value and the hassle avoided;
explicitly mention past purchases\\
\end{mmlongtable}

\subsection{Group 2 --- build relationships}

\begin{mmlongtable}{Group 2: relationship e-mails}{tab:group2}%
{@{}P{0.15\textwidth}P{0.28\textwidth}P{0.50\textwidth}@{}}
\rowcolor{ocre}\thd{Type} & \thd{Goal} & \thd{Rules}\\
\midrule
\endfirsthead
\rowcolor{ocre}\thd{Type} & \thd{Goal} & \thd{Rules}\\
\midrule
\endhead
\rlab{Newsletter} & Improve brand awareness & Think ``mini-newspaper''; break copy
into digestible sections; use simple layouts; always include contact
information\\
\rowcolor{tablealt}
\rlab{Welcome} & Establish a relationship with new opt-ins & Use a warm,
conversational tone --- a ``virtual handshake''; offer a reward such as 10\% off;
remind them of the benefits of joining\\
\rlab{Educational} & Build trust through industry knowledge & Offer relevant,
bite-sized information linking out to full blog posts; include a subtle
mini-promotion\\
\rowcolor{tablealt}
\rlab{Product Advice} & Establish authority & Solve customer problems; teach
product features; emphasise customer service; proofread meticulously to protect
credibility\\
\end{mmlongtable}

\subsection{Group 3 --- gather intelligence}

\begin{mmlongtable}{Group 3: intelligence e-mails}{tab:group3}%
{@{}P{0.15\textwidth}P{0.28\textwidth}P{0.50\textwidth}@{}}
\rowcolor{ocre}\thd{Type} & \thd{Goal} & \thd{Rules}\\
\midrule
\endfirsthead
\rowcolor{ocre}\thd{Type} & \thd{Goal} & \thd{Rules}\\
\midrule
\endhead
\rlab{Testimonials} & Reinforce business value through social proof & Use a sleek
design to frame quotes; pair text with powerful images --- customer photographs
or product shots; offer a link for readers to leave their own feedback\\
\rowcolor{tablealt}
\rlab{Surveys} & Collect data to improve customer experience & Clearly explain
the incentive, such as a prize draw; explicitly address why the data is being
collected; provide a highly visible, clickable access button\\
\end{mmlongtable}

\section{Transactional Against Promotional E-mail}

\begin{keybox}[High Yield]
``Describe the difference between transactional emails and promotional emails.
Provide examples of each and explain when to use them.''
\pyq{SEE, Oct/Nov 2023, Q6a --- 8 M}
\end{keybox}

\begin{mmlongtable}{Transactional against promotional e-mail}{tab:transvspromo}%
{@{}P{0.14\textwidth}P{0.40\textwidth}P{0.39\textwidth}@{}}
\rowcolor{ocre}\thd{Basis} & \thd{Transactional E-mail} & \thd{Promotional E-mail}\\
\midrule
\endfirsthead
\rowcolor{ocre}\thd{Basis} & \thd{Transactional E-mail} & \thd{Promotional E-mail}\\
\midrule
\endhead
\rlab{Trigger} & A specific action or event performed by the individual user &
A marketing calendar or campaign decision\\
\rowcolor{tablealt}
\rlab{Purpose} & To deliver information the user expects and needs & To generate
awareness, demand or sales\\
\rlab{Consent required} & Implied by the transaction itself & Explicit opt-in
required\\
\rowcolor{tablealt}
\rlab{Volume} & One-to-one, automated & One-to-many, batched or segmented\\
\rlab{Open rate} & Very high, often 40--60\% or more, because it is expected &
Much lower, typically 15--25\%\\
\rowcolor{tablealt}
\rlab{Examples} & Order confirmation, shipping and delivery notification,
password reset, one-time password, invoice or receipt, account creation,
appointment reminder & Sale announcement, newsletter, new-product launch,
seasonal offer, webinar invitation, re-engagement campaign\\
\rlab{When to use} & Immediately after the triggering event --- speed and clarity
matter more than design & At planned intervals against a content calendar,
targeted to segments\\
\end{mmlongtable}

\begin{keybox}[A Practical Point Worth Marks]
Because transactional e-mails are opened far more often, a small, tasteful
promotional element --- a related-product strip or a loyalty reminder ---
\emph{inside} them often outperforms an entire promotional campaign.

\smallskip
But the transactional content must remain dominant, or the message loses its
legal status as transactional and becomes subject to opt-in rules.
\end{keybox}

\begin{summarybox}
E-mail marketing delivers commercial and relationship messages to a
permission-based list in order to acquire, engage, convert and retain. It is the
only major owned channel, which is why it has the highest conversion and
retention rate. E-mail 2.0 replaces one-way blasts and static auto-responders
with behavioural triggers, segmentation and lifecycle automation. Its advantages
are cost, speed, tracking, personalisation, reach and ownership; its liabilities
are legal risk, engagement decay, rendering inconsistency, hidden cost,
deliverability management and inbox competition. Opt-in differs from spam in
acquisition, expectation, brand impact and list quality; double opt-in produces a
smaller and cleaner list. An effective e-mail keeps to one topic, a simple design
and one clear call to action. The nine templates group into revenue,
relationship and intelligence. Transactional e-mail is user-triggered, expected
and opened two to three times as often as promotional e-mail, which is
calendar-driven and requires explicit consent.
\end{summarybox}

%%%%%%%%%%%%%%%%%%%%%%%%%%%%%%%%%%%%%%%%%%%%%%%%%%%%%%%%%%%%%%%%%%%%%%%%%%%%%%
\chapter{The E-mail Plan}
%%%%%%%%%%%%%%%%%%%%%%%%%%%%%%%%%%%%%%%%%%%%%%%%%%%%%%%%%%%%%%%%%%%%%%%%%%%%%%

\syllabusstrip{E-mail Plan.}

\begin{learningoutcomes}
\begin{itemize}
\item The seven-node e-mail marketing plan, and why it is a loop.
\item A worked plan for a stated scenario.
\item The eight-step guide to segmentation and personalisation.
\end{itemize}
\end{learningoutcomes}

\section{The Seven Nodes}

\begin{keybox}[The Framework for Every ``Design an E-mail Plan'' Question]
There have been several such questions, and each is answered by walking these
seven nodes against the specific business given.
\pyq{CIE-II, 26 May 2026 --- 10 M}
\end{keybox}

\begin{figure}[htbp]
\centering
\begin{tikzpicture}[font=\sffamily\footnotesize,
    nd/.style={draw=ocre,line width=0.9pt,fill=ocrepale,circle,
               minimum size=1.72cm,align=center,
               font=\sffamily\bfseries\scriptsize,text=slate,inner sep=1pt}]
  \node[nd] (n1) at (0,0)     {1\\ choose\\ software};
  \node[nd] (n2) at (1.95,0)  {2\\ build\\ the list};
  \node[nd] (n3) at (3.90,0)  {3\\ capture\\ forms};
  \node[nd] (n4) at (5.85,0)  {4\\ define\\ goals};
  \node[nd] (n5) at (7.80,0)  {5\\ auto-\\ responders};
  \node[nd] (n6) at (9.75,0)  {6\\ add\\ triggers};
  \node[nd,draw=greenok,fill=greenpale] (n7) at (11.70,0) {7\\ monitor\\ metrics};
  \foreach \a/\b in {n1/n2,n2/n3,n3/n4,n4/n5,n5/n6,n6/n7}{
    \draw[-{Latex[length=2mm]},ocre,line width=0.9pt] (\a)--(\b);}
  \draw[-{Latex[length=2.2mm]},greenok,line width=0.9pt,dashed]
    (n7.south) -- ++(0,-1.05) -| (n1.south);
  \node[anchor=north,text=greenok,font=\sffamily\scriptsize\itshape]
    at (5.85,-1.1)
    {the metrics tell you whether the software, list, goals, sequences and
     triggers were right};
\end{tikzpicture}
\caption[The seven-node e-mail marketing plan]%
        {The seven nodes form a loop, not a line. Node~7 feeds back into Node~1.}
\label{fig:sevennodes}
\end{figure}

\begin{mmlongtable}{The seven nodes in detail}{tab:sevennodes}%
{@{}P{0.20\textwidth}P{0.73\textwidth}@{}}
\rowcolor{ocre}\thd{Node} & \thd{Action}\\
\midrule
\endfirsthead
\rowcolor{ocre}\thd{Node} & \thd{Action}\\
\midrule
\endhead
\rlab{1. Choose software} & Select a CRM or e-mail service provider based on the
capabilities required --- a simple database sender against full behavioural
campaign automation\\
\rowcolor{tablealt}
\rlab{2. Build the list} & Start capturing addresses; commit to quality over
quantity; limit sending to a maximum of about two e-mails per month initially\\
\rlab{3. Capture forms} & Deploy sign-up forms on the website and blog to gather
names, e-mail addresses and phone numbers\\
\rowcolor{tablealt}
\rlab{4. Define goals} & Determine the mission --- deepen relationships, educate
the audience, or drive sales cycles\\
\rlab{5. Set auto-responders} & Build a sequence of at least six educational,
non-salesy e-mails released automatically to new subscribers\\
\rowcolor{tablealt}
\rlab{6. Add triggers} & Automate new sales cycles based on specific user
actions --- for example clicking a particular link, or abandoning a cart\\
\rlab{7. Monitor metrics} & Review reporting monthly to test subject lines and
improve click-through rate --- and feed the insight back into Node~1\\
\end{mmlongtable}

\section{Worked Plan --- a New Online Business Targeting College Students}

\begin{exambox}[Official Scheme Answer --- 10 Marks]
\pyq{CIE-II, 26 May 2026 --- 10 M}
\textbf{Objective:} increase awareness, engagement and sales among college
students. \textbf{Target audience:} college students aged 18--25.

\begin{enumerate}
\item \sfb{Build the e-mail list} --- collect e-mails through website sign-ups;
      offer discounts and free resources.
\item \sfb{Create attractive content} --- use youthful, creative designs;
      include offers, contests and product updates.
\item \sfb{Segmentation} --- segment users based on interests and purchase
      behaviour.
\item \sfb{Personalisation} --- use customer names and personalised
      recommendations.
\item \sfb{Campaign schedule} --- send weekly newsletters and promotional
      e-mails.
\item \sfb{Mobile optimisation} --- ensure e-mails are mobile-friendly.
\item \sfb{Performance tracking} --- track open rate, click-through rate\ix{click-through rate},
      conversion rate and unsubscribe rate.
\end{enumerate}
\end{exambox}

\section{Segmentation and Personalisation, Step by Step}

\begin{keybox}[High Yield]
``Describe the importance of audience segmentation and personalization in email
marketing. Provide a step-by-step guide on how to segment an audience effectively
and personalize email content based on segmentation criteria.''
\pyq{CIE-II, 24 Jul 2024 --- 10 M}
\end{keybox}

\begin{enumerate}
\item \sfb{Collect the right data at capture} --- ask only for what you will use,
      then enrich progressively over time.
\item \sfb{Choose segmentation criteria} --- demographics (age, location,
      gender); firmographics for B2B (industry, company size, role); behaviour
      (pages viewed, products browsed, past purchases, cart abandonment);
      engagement level (active, lapsing, dormant); lifecycle stage (new
      subscriber, first-time buyer, repeat customer, advocate); and preference
      (the topics and frequency the subscriber selected).
\item \sfb{Build the segments} in the e-mail service provider as \emph{dynamic
      rules}, so that membership updates automatically.
\item \sfb{Design content per segment} --- a different subject line, offer,
      imagery and call to action for each.
\item \sfb{Layer personalisation inside the segment} --- merge tags for the name,
      dynamic product blocks based on browsing history, location-specific offers,
      and send-time optimisation.
\item \sfb{Automate lifecycle triggers} --- welcome series, browse abandonment,
      cart abandonment, post-purchase, replenishment, win-back.
\item \sfb{Test} --- A/B test subject lines and offers \emph{within} each
      segment, since what works for one will not work for another.
\item \sfb{Measure per segment}, and prune or re-engage the segments that
      consistently underperform.
\end{enumerate}

\begin{examplebox}[A Fashion Retailer Segmenting by Past Category]
Subscribers who bought running shoes receive a new-arrivals e-mail featuring
running gear with the subject line ``New for your next 10K'', while subscribers
who bought formal wear receive an entirely different creative.

\smallskip
The same send, two segments --- and typically double the click-through rate of a
single generic blast.
\end{examplebox}

\begin{summarybox}
The e-mail plan has seven nodes --- choose software, build the list, deploy
capture forms, define goals, set auto-responders, add behavioural triggers, and
monitor metrics --- and Node~7 loops back to Node~1. Applied to a scenario, the
plan becomes a specific answer: list building, creative, segmentation,
personalisation, schedule, mobile optimisation and tracking. Segmentation and
personalisation run in eight steps: collect the right data, choose criteria,
build dynamic segments, design per-segment content, layer personalisation inside
the segment, automate lifecycle triggers, test within segments, and measure per
segment.
\end{summarybox}

%%%%%%%%%%%%%%%%%%%%%%%%%%%%%%%%%%%%%%%%%%%%%%%%%%%%%%%%%%%%%%%%%%%%%%%%%%%%%%
\chapter{E-mail Marketing Campaign Analysis}
%%%%%%%%%%%%%%%%%%%%%%%%%%%%%%%%%%%%%%%%%%%%%%%%%%%%%%%%%%%%%%%%%%%%%%%%%%%%%%

\syllabusstrip{E-mail Marketing Campaign Analysis.}

\begin{learningoutcomes}
\begin{itemize}
\item What campaign analysis means.
\item The six-stage measurement funnel, with formulas.
\item How to read each KPI, what a poor reading means, and the fix.
\item The challenges of running e-mail campaigns, with paired solutions.
\item Two concrete strategies for raising open and click rates.
\end{itemize}
\end{learningoutcomes}

\section{What Campaign Analysis Means}

\begin{exambox}[Model Answer --- 2 Marks]
\pyq{CIE-II, 26 May 2026 --- 2 M}
Campaign analysis in e-mail marketing refers to evaluating the performance of
e-mail campaigns using metrics such as open rate, click-through rate, conversion
rate, bounce\ix{bounce} rate and unsubscribe rate, in order to determine campaign
effectiveness and improve future campaigns.
\end{exambox}

\section{The Measurement Funnel}

The funnel runs from the top of the list to the bottom line. Each stage narrows,
and each has its own formula and diagnostic meaning.

\begin{figure}[htbp]
\centering
\begin{tikzpicture}[font=\sffamily\footnotesize]
  % funnel bands, narrowing downward
  \fill[ocre!45]  (-4.6,0)     -- (4.6,0)     -- (4.0,-1.0)  -- (-4.0,-1.0)  -- cycle;
  \fill[ocre!36]  (-4.0,-1.0)  -- (4.0,-1.0)  -- (3.4,-2.0)  -- (-3.4,-2.0)  -- cycle;
  \fill[ocre!28]  (-3.4,-2.0)  -- (3.4,-2.0)  -- (2.8,-3.0)  -- (-2.8,-3.0)  -- cycle;
  \fill[ocre!21]  (-2.8,-3.0)  -- (2.8,-3.0)  -- (2.2,-4.0)  -- (-2.2,-4.0)  -- cycle;
  \fill[ocre!14]  (-2.2,-4.0)  -- (2.2,-4.0)  -- (1.6,-5.0)  -- (-1.6,-5.0)  -- cycle;
  \fill[greenok!30] (-1.6,-5.0) -- (1.6,-5.0) -- (1.0,-6.0)  -- (-1.0,-6.0)  -- cycle;

  \node[text=slate,font=\sffamily\bfseries\scriptsize] at (0,-0.5)  {1.\ \ SENT E-MAILS};
  \node[text=slate,font=\sffamily\bfseries\scriptsize] at (0,-1.5)  {2.\ \ DELIVERY (NON-BOUNCE)};
  \node[text=slate,font=\sffamily\bfseries\scriptsize] at (0,-2.5)  {3.\ \ OPEN RATE};
  \node[text=slate,font=\sffamily\bfseries\scriptsize] at (0,-3.5)  {4.\ \ CLICK-THROUGH RATE};
  \node[text=slate,font=\sffamily\bfseries\scriptsize] at (0,-4.5)  {5.\ \ CONVERSION RATE};
  \node[text=greenok,font=\sffamily\bfseries\scriptsize] at (0,-5.5) {6.\ \ REVENUE};

  \node[anchor=west,text=slatelight,font=\sffamily\scriptsize,align=left] at (4.8,-0.5)
    {list growth rate};
  \node[anchor=west,text=slatelight,font=\sffamily\scriptsize,align=left] at (4.8,-1.5)
    {list hygiene, sender reputation};
  \node[anchor=west,text=slatelight,font=\sffamily\scriptsize,align=left] at (4.8,-2.5)
    {subject line, sender name, timing};
  \node[anchor=west,text=slatelight,font=\sffamily\scriptsize,align=left] at (4.8,-3.5)
    {content relevance, CTA clarity};
  \node[anchor=west,text=slatelight,font=\sffamily\scriptsize,align=left] at (4.8,-4.5)
    {offer quality, landing page};
  \node[anchor=west,text=greenok,font=\sffamily\scriptsize,align=left] at (4.8,-5.5)
    {revenue per message};

  \node[anchor=east,rotate=90,text=ocredark,font=\sffamily\bfseries\scriptsize]
    at (-5.0,-3.0) {THE MARKETING CONTROL ROOM};
\end{tikzpicture}
\caption[The e-mail measurement funnel]%
        {The measurement funnel. Each stage has a distinct diagnosis, so a fall
         in revenue can be traced to exactly one narrowing --- which is why the
         funnel must be read as a system rather than as six separate numbers.}
\label{fig:emailfunnel}
\end{figure}

\begin{mmlongtable}{The six-stage measurement funnel}{tab:funnel}%
{@{}P{0.18\textwidth}P{0.42\textwidth}P{0.33\textwidth}@{}}
\rowcolor{ocre}\thd{Stage} & \thd{Definition or formula} & \thd{What it diagnoses}\\
\midrule
\endfirsthead
\rowcolor{ocre}\thd{Stage} & \thd{Definition or formula} & \thd{What it diagnoses}\\
\midrule
\endhead
\rlab{1. Sent e-mails} & The starting list volume. Metric: list growth rate &
Whether the audience asset is growing or decaying\\
\rowcolor{tablealt}
\rlab{2. Delivery (non-bounce rate)} & $100\% - (\text{bounced} \div
\text{total sent})$. Accounts for invalid addresses or full firewalls &
List hygiene and sender reputation\\
\rlab{3. Open rate} & Tracked when images render on the server or links are
clicked. Calculated from the non-bounce pool & The strength of the subject line,
sender name and send timing\\
\rowcolor{tablealt}
\rlab{4. Click-through rate} & The percentage of opens that clicked a link. Used
to track interests and A/B test variations & The relevance of the content and the
clarity of the call to action\\
\rlab{5. Conversion rate} & The percentage of clicks that completed the ultimate
goal --- purchasing, filling a form & The quality of the offer and the landing
page\\
\rowcolor{tablealt}
\rlab{6. Revenue per message} & $\text{total revenue} \div \text{number of sent
e-mails}$ & The single bottom-line measure of the whole programme\\
\end{mmlongtable}

\begin{exambox}[Define Click-Through Rate --- 2 Marks]
\pyq{CIE-II, 26 May 2026 --- 2 M}
Click-through rate is the percentage of recipients who click on links in an
e-mail compared to the total number of \emph{delivered} e-mails.
\end{exambox}

\begin{warnbox}[Two Denominators, Two Names]
The lecture material defines CTR against \emph{opens}; the scheme of valuation
defines it against \emph{delivered} e-mails. The industry distinguishes these as
\textbf{click-to-open rate} and \textbf{click-through rate} respectively.

\smallskip
State the formula you are using and either is defensible. Stating no formula is
what loses the mark.
\end{warnbox}

\section{Interpreting the KPIs}

\begin{keybox}[High Yield]
``Analyse the effectiveness of an E-mail Marketing campaign using KPIs such as
open rate, click-through rate, unsubscribe rate, and conversion rate.''
\pyq{CIE-II, 26 May 2026 --- 10 M}
Anchor the answer in the six-stage funnel above, so that the KPIs appear as a
connected system rather than a list.
\end{keybox}

\begin{mmlongtable}{Reading the e-mail KPIs}{tab:emailkpi}%
{@{}P{0.157\textwidth}P{0.235\textwidth}P{0.255\textwidth}P{0.265\textwidth}@{}}
\rowcolor{ocre}\thd{KPI} & \thd{What it measures} & \thd{What a poor reading means} & \thd{Fix}\\
\midrule
\endfirsthead
\rowcolor{ocre}\thd{KPI} & \thd{What it measures} & \thd{What a poor reading means} & \thd{Fix}\\
\midrule
\endhead
\rlab{Open rate} & How many users opened the e-mail & Weak subject line, wrong
send time, poor sender reputation & Test subject lines; clean the list;
authenticate the domain\\
\rowcolor{tablealt}
\rlab{Click-through rate} & How many users clicked links in the e-mail & Content
not engaging, or the call to action unclear & Improve the offer; make one CTA
dominant; segment better\\
\rlab{Unsubscribe rate} & Users opting out of e-mails & Content is irrelevant, or
frequency is too high & Send relevant content; offer a preference centre instead
of only ``unsubscribe''\\
\rowcolor{tablealt}
\rlab{Conversion rate} & Users completing the desired action & The offer or
landing page is failing, not the e-mail & Fix message match on the landing page;
reduce form friction\\
\rlab{Bounce rate} & Undelivered e-mails & Stale or purchased list; server blocks
& Remove hard bounces immediately; use double opt-in\\
\rowcolor{tablealt}
\rlab{Spam complaint rate} & Recipients marking the e-mail as junk & Consent is
unclear, or expectations were not set & Set expectations at sign-up; make
unsubscribe obvious\\
\end{mmlongtable}

\begin{keybox}[The Diagnostic Logic]
Read the funnel downward and the failure locates itself. Clicks but no
conversions is \emph{not} an e-mail problem --- the e-mail did its job and handed
the visitor to a page that failed. Opens but no clicks \emph{is} an e-mail
problem: the subject line promised something the body did not deliver.
\end{keybox}

\section{Challenges and Solutions}

\begin{exambox}[Official Scheme Answer --- 10 Marks]
\pyq{CIE-II, 26 May 2026 --- 10 M}
\end{exambox}

\begin{mmlongtable}{E-mail marketing challenges and solutions}{tab:challenges}%
{@{}P{0.38\textwidth}P{0.55\textwidth}@{}}
\rowcolor{ocre}\thd{Challenge} & \thd{Solution}\\
\midrule
\endfirsthead
\rowcolor{ocre}\thd{Challenge} & \thd{Solution}\\
\midrule
\endhead
\rlab{Low open rates} & Use attractive subject lines\\
\rowcolor{tablealt}
\rlab{Spam filtering issues} & Avoid spam words and maintain sender reputation\\
\rlab{High unsubscribe rates} & Send relevant and useful content\\
\rowcolor{tablealt}
\rlab{Poor e-mail design} & Use responsive e-mail templates\\
\rlab{Lack of personalisation} & Segment the audience and personalise e-mails\\
\rowcolor{tablealt}
\rlab{Inactive subscribers} & Run re-engagement campaigns\\
\end{mmlongtable}

\noindent Additional challenges worth naming: deliverability and inbox
placement; \keytermas{list decay}{list decay} of roughly 20--25\% per year;
rendering inconsistency across clients; privacy changes such as mail privacy
protection that inflate open rates and make them unreliable; and a rising
regulatory burden.

\section{Raising Open and Click Rates}

\begin{exambox}[Two Strategies --- 2 Marks, for a Fashion Brand's Summer Sale]
\pyq{CIE-II Quiz, 24 Jul 2024 --- 2 M}
\begin{itemize}
\item \sfb{To increase open rate:} write a specific, urgency-bearing,
      personalised subject line and A/B test it --- for example ``Priya, your
      summer edit is 40\% off --- today only'' --- and optimise the send time to
      the segment's historical open window.
\item \sfb{To increase click-through rate:} use a single dominant, benefit-led
      call-to-action button above the fold, supported by a strong hero product
      image and segmented product recommendations based on the recipient's
      browsing history.
\end{itemize}
\end{exambox}

\begin{summarybox}
Campaign analysis evaluates e-mail performance on open rate, click-through rate,
conversion rate, bounce rate and unsubscribe rate in order to improve future
campaigns. The measurement funnel has six stages --- sent, delivered, opened,
clicked, converted, revenue --- each with a formula and a distinct diagnosis, so
a fall in revenue traces to exactly one narrowing. Each KPI has a characteristic
failure and a matching fix, and the key diagnostic distinction is that clicks
without conversions indicts the landing page while opens without clicks indicts
the e-mail. Six standard challenges pair with six solutions, and to those should
be added deliverability, list decay of 20--25\% a year, rendering inconsistency
and privacy changes that make open rate unreliable.
\end{summarybox}

%%%%%%%%%%%%%%%%%%%%%%%%%%%%%%%%%%%%%%%%%%%%%%%%%%%%%%%%%%%%%%%%%%%%%%%%%%%%%%
\chapter{Measuring Conversions and Keeping Up}
%%%%%%%%%%%%%%%%%%%%%%%%%%%%%%%%%%%%%%%%%%%%%%%%%%%%%%%%%%%%%%%%%%%%%%%%%%%%%%

\syllabusstrip{Measuring Conversions \& keeping up.}

\begin{learningoutcomes}
\begin{itemize}
\item How to track conversions that happen offline.
\item How to integrate web analytics with e-mail marketing.
\item The E-mail 2.0 technology stack.
\item Responsive e-mail design --- the case, the method and the critical view.
\item The four pillars of the e-mail engine.
\end{itemize}
\end{learningoutcomes}

\section{Bridging the Gap --- Offline Tracking}

When open rates and click-through rates are not enough to track overall marketing
success, e-mail requires physical participation mechanisms.

\begin{mmlongtable}{Tracking conversions that happen offline}{tab:offline}%
{@{}P{0.22\textwidth}P{0.71\textwidth}@{}}
\rowcolor{ocre}\thd{Mechanism} & \thd{How it is tracked}\\
\midrule
\endfirsthead
\rowcolor{ocre}\thd{Mechanism} & \thd{How it is tracked}\\
\midrule
\endhead
\rlab{In-store purchases} & By requesting the customer's e-mail address at the
physical point of sale and matching it against the campaign list\\
\rowcolor{tablealt}
\rlab{Phone calls} & By providing dedicated phone numbers exclusive to the
e-mail campaign\\
\rlab{Event attendance} & By issuing special entrance codes that can only be
accessed and viewed from within the e-mail itself\\
\end{mmlongtable}

\begin{keybox}[The Digital Equivalents of the Same Principle]
Unique promotion codes per campaign; UTM-tagged links feeding web analytics; and
printable or scannable QR coupons that carry a campaign identifier. In every case
the mechanism is the same: give this campaign, and only this campaign, a token
that can be counted later.
\end{keybox}

\section{Integrating Web Analytics and E-mail Marketing}

\begin{exambox}[Official Scheme Answer --- 10 Marks]
\pyq{CIE-II, 26 May 2026 --- 10 M}
\begin{enumerate}
\item \sfb{Collect customer data} --- use web analytics tools to track user
      behaviour, page visits and purchase history.
\item \sfb{Audience segmentation} --- segment users based on interests and
      browsing behaviour.
\item \sfb{Personalised e-mail campaigns} --- send customised e-mails with
      relevant product recommendations.
\item \sfb{Track campaign performance} --- use open rate, click-through rate,
      conversion rate and bounce rate.
\item \sfb{Retargeting} --- send follow-up e-mails to abandoned-cart users.
\item \sfb{Continuous optimisation} --- analyse campaign data and improve content
      and timing.
\end{enumerate}

\smallskip
\textbf{Benefits:} better customer engagement; higher conversion rates; improved
customer retention\ix{customer retention}; increased sales and revenue.
\end{exambox}

\section{The E-mail 2.0 Technology Stack}

\begin{mmlongtable}{The e-mail technology stack}{tab:emailstack}%
{@{}P{0.24\textwidth}P{0.69\textwidth}@{}}
\rowcolor{ocre}\thd{Category} & \thd{Tools}\\
\midrule
\endfirsthead
\rowcolor{ocre}\thd{Category} & \thd{Tools}\\
\midrule
\endhead
\rlab{CRM and all-in-one automation} & HubSpot (automated workflows),
Infusionsoft (purchase behaviour and task management), ActiveCampaign (lifecycle
automation), Contactually (social and inbox vetting)\\
\rowcolor{tablealt}
\rlab{Drip campaigns and nurturing} & Drip (lightweight workflow builder),
ConvertKit (automation rules for publishers), Emma (personalisation via CRM
sync), Campaign Monitor (interactive analytics), Sendloop\\
\rlab{Creation and standard sending} & MailChimp (abandoned carts, groups),
GetResponse (calendar-integrated auto-responders), AWeber (new-subscriber
catch-up), Reach Mail (simultaneous comparison of five campaigns)\\
\rowcolor{tablealt}
\rlab{Specialised utility} & Litmus (one-click render testing across around 40
clients), BombBomb (video-powered e-mail), Gumroad (digital product updates),
Intercom (in-app conversations)\\
\rlab{Sales tracking plugins} & GetNotify (invisible open tracking), ToutApp
(playbooks and CRM sync), Yesware (prescriptive analytics, team tracking),
Clearslide (web engagement tracking)\\
\end{mmlongtable}

\section{Responsive E-mail Design}

\begin{keybox}[High Yield]
``Evaluate the effectiveness of responsive email design in maintaining a
competitive edge in digital marketing.''
\pyq{SEE, Jun/Jul 2025, Q6a --- 8 M}
The verb is \emph{evaluate}, so the answer must reach a judgement rather than
merely describe the technique.
\end{keybox}

\subsection{The case for it}

The majority of e-mail opens now occur on mobile, and a non-responsive e-mail
requiring pinch-and-zoom is typically deleted within seconds. Responsive design\ix{responsive design}
therefore protects open-to-click conversion, reduces unsubscribes, improves
accessibility, and preserves brand perception. It also improves deliverability
indirectly, because engagement signals feed inbox placement algorithms.

\subsection{How it is achieved}

\begin{itemize}
\item Fluid single-column layouts and media queries.
\item Scalable images.
\item Body text of at least 14 to 16 pixels.
\item Tap targets of at least 44 by 44 pixels.
\item ``Bulletproof'' call-to-action buttons built in HTML rather than as images.
\item Alt text, for image-blocked clients.
\end{itemize}

\subsection{The critical view}

\begin{warnbox}[Responsiveness Is Table Stakes, Not a Differentiator]
Every competent competitor has responsive e-mail, so it prevents disadvantage
more than it creates advantage.

\smallskip
Genuine competitive edge comes from \textbf{segmentation, personalisation,
behavioural triggers and offer quality}. Responsive design is the
\emph{precondition} that allows those to work.

\smallskip
It also adds testing overhead across dozens of clients --- hence render-testing
tools --- and dark-mode rendering remains inconsistent.
\end{warnbox}

\section{Operating the Engine --- Four Pillars}

\begin{figure}[htbp]
\centering
% The pillar shafts are drawn as fixed rectangles and the text is placed as
% separately anchored nodes inside them, so no amount of text can push a shaft
% through the capital band above it.
\begin{tikzpicture}[font=\sffamily\footnotesize,
    hd/.style={anchor=north,align=center,text width=2.40cm,
               font=\sffamily\bfseries\scriptsize,text=ocredark},
    sb/.style={anchor=north,align=center,text width=2.40cm,
               font=\sffamily\scriptsize\itshape,text=slatelight},
    bd/.style={anchor=north,align=center,text width=2.40cm,
               font=\sffamily\scriptsize,text=slate}]

  % capital and base
  \fill[ocre]  (-0.35,4.80) rectangle (11.50,5.10);
  \fill[slate] (-0.35,-0.40) rectangle (11.50,-0.08);
  \node[anchor=north,text=slate,font=\sffamily\bfseries\scriptsize]
    at (5.58,-0.55) {A TARGETED, MEASURABLE, HIGH-ROI ECOSYSTEM};

  % ---- pillar 1
  \fill[ocrepale] (0.0,0) rectangle (2.6,4.75);
  \draw[ocre,line width=1pt] (0.0,0) rectangle (2.6,4.75);
  \node[hd] at (1.3,4.62) {PILLAR 1\\ THE FUEL};
  \node[sb] at (1.3,3.85) {the opt-in list};
  \node[bd] at (1.3,3.42) {micro-focused, permission-based audiences that
     eliminate legal risk and guarantee engagement};

  % ---- pillar 2
  \fill[ocrepale] (2.85,0) rectangle (5.45,4.75);
  \draw[ocre,line width=1pt] (2.85,0) rectangle (5.45,4.75);
  \node[hd] at (4.15,4.62) {PILLAR 2\\ THE PAYLOAD};
  \node[sb] at (4.15,3.85) {content arsenal};
  \node[bd] at (4.15,3.42) {deploy the exact right template: promotional,
     educational, survey};

  % ---- pillar 3
  \fill[ocrepale] (5.70,0) rectangle (8.30,4.75);
  \draw[ocre,line width=1pt] (5.70,0) rectangle (8.30,4.75);
  \node[hd] at (7.0,4.62) {PILLAR 3\\ THE ENGINE};
  \node[sb] at (7.0,3.85) {triggers and tech};
  \node[bd] at (7.0,3.42) {move from static blasts to dynamic behavioural
     lifecycle automation};

  % ---- pillar 4
  \fill[ocrepale] (8.55,0) rectangle (11.15,4.75);
  \draw[ocre,line width=1pt] (8.55,0) rectangle (11.15,4.75);
  \node[hd] at (9.85,4.62) {PILLAR 4\\ THE DASHBOARD};
  \node[sb] at (9.85,3.85) {the metrics};
  \node[bd] at (9.85,3.42) {obsess over the funnel, from non-bounce delivery to
     offline revenue};
\end{tikzpicture}
\caption[The four pillars of the e-mail engine]%
        {The four pillars. Remove any one and the engine stalls: a perfect
         template sent to a purchased list is illegal, and a perfectly automated
         programme without a dashboard cannot be improved.}
\label{fig:fourpillars}
\end{figure}

\begin{mmlongtable}{The four pillars}{tab:pillars}%
{@{}P{0.28\textwidth}P{0.65\textwidth}@{}}
\rowcolor{ocre}\thd{Pillar} & \thd{What it means}\\
\midrule
\endfirsthead
\rowcolor{ocre}\thd{Pillar} & \thd{What it means}\\
\midrule
\endhead
\rlab{1 --- The Fuel (opt-in list)} & Cultivate micro-focused, permission-based
audiences, to eliminate the legal risk and guarantee engagement\\
\rowcolor{tablealt}
\rlab{2 --- The Payload (content arsenal)} & Deploy the exact right template ---
promotional to drive revenue, educational to build trust, surveys to gather
intelligence\\
\rlab{3 --- The Engine (triggers and tech)} & Use the technology stack to move
from static blasts to dynamic, behavioural, lifecycle automation\\
\rowcolor{tablealt}
\rlab{4 --- The Dashboard (metrics)} & Obsess over the measurement funnel ---
from non-bounce delivery down to total offline revenue generation\\
\end{mmlongtable}

\begin{keybox}[The Closing Idea]
E-mail marketing is no longer a megaphone. It is a targeted, measurable, high-ROI
ecosystem.
\end{keybox}

\begin{summarybox}
Conversions that happen offline are tracked by giving each campaign a countable
token: an e-mail address captured at the point of sale, a dedicated phone number,
an e-mail-only entrance code, or their digital equivalents in promotion codes,
UTM tags and QR coupons. Web analytics and e-mail integrate in six steps ---
collect behavioural data, segment, personalise, track performance, retarget
abandoners, and optimise continuously --- yielding better engagement, higher
conversion, improved retention and increased revenue. The technology stack spans
CRM automation, drip nurturing, standard sending, specialised utilities and sales
tracking. Responsive design protects open-to-click conversion and is now a
precondition rather than a differentiator; the real edge lies in segmentation,
personalisation, triggers and offer quality. Four pillars operate the engine:
the opt-in list, the content arsenal, the trigger technology and the metrics
dashboard.
\end{summarybox}
